
\subsection{Idée 1}

Suivant la stratégie développée par C.Duval \cite[Théorème 2.2]{duval2013density}, nous allons chercher à établir une inégalité minimax pour la super--expérience $\Tilde Z$ pour lequel nous observons les trajectoires de manière continue, \ie que nous supposons que l'on a observé les différents temps de sauts jusqu'au instant $t_1$ et $t_2$.  \\
L'idée principale, consiste à dire que si nous sommes capable d'obtenir une borne inférieure pour la super--expérience, alors une borne inférieure sur le vrai processus $Z$ ne peut qu'être supérieure. \\

Soient $\PP_{T,0}$ la loi de $\Tilde Z$ et $\PP_{T,t_1, t_2}$ la loi des observations $Z = (Z_{t_1}^j, Z_{t_2}^j)_{j = 1\cdots J}$. 

\begin{align*}
    \inf_{\hat f} \sup_{f \in \Fcal}  \Big( \EE_{\PP_{T,t_1, t_2}} \Big[  \| \hat f - f \|_p^p  \Big] \Big)^{1/p}
    \geq 
    \inf_{\tilde f} \sup_{f \in \Fcal}  \Big( \EE_{\PP_{T,0}} \Big[  \| \tilde f - f \|_p^p  \Big] \Big)^{1/p}    
    \eqfinp
\end{align*}

\subsubsection{Borne inférieure pour la super expérience}

Donnons nous une mesure sigma-finie sur $\Fcal$.  Alors
\begin{align*}
    \sup_{f \in \Fcal}  \EE_{\PP_{T,0}} \Big[  \| \tilde f - f \|_p^p  \Big| N_{t_1} = k_1; N_{t_2} = k_2 \Big]  
    &\geq
   \int_{\Fcal}  \EE_{\PP_{T,0}} \Big[  \| \tilde f - f \|_p^p  \Big| N_{t_1} = k_1; N_{t_2} = k_2 \Big]  \diff \mu(f)
\end{align*}

Mais, par Fubini
\begin{align*}
   \sum_{k_1 = 0}^{+\infty} \sum_{k_2 = 0}^{+\infty}
   \nu(k_1, k_2) \int_{\Fcal}  \EE_{\PP_{T,0}} \Big[  \| \tilde f - f \|_p^p  \Big| N_{t_1} = k_1; N_{t_2} = k_2 \Big]  \diff \mu(f)
   =
   \int_{\Fcal}  \EE_{\PP_{T,0}} \Big[  \| \tilde f - f \|_p^p \Big]  \diff \mu(f) 
\end{align*}

\begin{question} Dans la thèse de Duval, 
\cite{duval2013density}, elle établit la borne inférieure sur son processus en utilisant un résultat de Donoho et al. \cite[Théorème 2]{donoho1996density}. Le point important de ce résultat c'est qu'on regarde des réalisations indépendantes d'un processus, identiquement distribuées. 
    \begin{itemize}
        \item Le truc, c'est que dans mon cas je regarde le processus en deux instants $t_1$ et $t_2$. J'ai donc une série d'observations $Z_{t_1}^j = \sum\limits_{k = 0}^{N_{t_1}^j} X_k^j + \varepsilon_{t_1}^j$ et $Z_{t_2}^j =  \sum\limits_{k = 0}^{N_{t_2}^j} X_k^j + \varepsilon_{t_2}^j$ 
        Pour avoir des observations indépendantes, je peux considérer que j'ai les observations 
        $(Z_{t_1}^j)_j$ et $(Z_{t_2}^j - Z_{t_1}^j)_j$ qui ne sont pas \iid. \\

        Dans le cas de Duval \cite{duval2013density}, c'est plus simple, car il n'y a pas de temps différents. Les incréments sont donc bien \iid. 
    \end{itemize}
\end{question}


\subsection{Preuve minimax : cas bizzare}


To control the integral, we must control $|e^{\lambda t_1 \varphi_1(x)} - e^{\lambda t_1 \varphi_0(x)}|$
. For that, we must control to part : 
\begin{itemize}
    \item If $|x| \to +\infty$. In that case, we know that $f_j$ and $f_0$ behaves as $x^{-2k}$.
    A asymptotic expansion ensures that
    \begin{align*}
        |e^{\lambda t_1 f_1(x)} - e^{\lambda t_1 f_0(x)}|
        = |
        \lambda t_1| \cdot |f_1(x) - f_0(x)| = O \Big(\frac{1}{x^{2k}} \Big). 
    \end{align*}
    \item If $|x| \to 0$. A characteristic function is always bounded by 1. It follows that 
    \begin{align*} 
    \big(\varphi_{1, Z_{t_1}}(x) - \varphi_{0, Z_{t_1}}(x)\big)^2
    \leq 4. \\
    \end{align*}
    \end{itemize}

$\blacktriangleright$ It remains to control the term
    \begin{align*}
     \int_{-\infty}^{+\infty} 
     \big(\varphi_{1, Z_{t_1}}(x)^{'k} - \varphi_{0, Z_{t_1}}(x)^{'k}\big)^2 \diff x 
     \eqfinp
    \end{align*}
    \begin{itemize}
        \item If $|x| \to +\infty$, the previous method allows us to obtain a bound of the form $O \big( \frac1{x^p}\big)$ where $p$ only depend on the regularity of $\varphi_\varepsilon$ et de $k$. 
        \item If $|x| \to 0$. Here, we can't use the fact that a characteristic function is bounded by $1$ because we have no information on the derivative of $\varphi_{1, Z_{t_1}}$. \\
        However, if $k = 1$. We have
            \begin{align*}
            \varphi'_{1, Z_{t_1}}(x)
            &=
            \lambda t_1 f_0'(x) \varphi_{1, Z_{t_1}}(x) + e^{-\lambda t_1 + \lambda t_1 f_0(x)} \varphi_\varepsilon'(x) \\
            &\leq 
            \lambda t_1 f_0'(x) + \varphi_\varepsilon'(x)
            \end{align*}
        Then, we need to add more assumption on $f_0$ and $\varphi_\varepsilon$ to be able to control our term. 
        \end{itemize}